\documentclass{beamer}
% Concise presentation for a 5-minute talk
\title{K-means: Implementation and Evaluation}
\author{Emilio Dalla Torre, Efrem Marzatico,\\Sebastiano Perni, Lorenzo Terzi}
\institute{Hands-on K-means}
\date{\today}

\begin{document}

\frame{\titlepage}

\begin{frame}\frametitle{Objectives}
    \begin{itemize}
        \item Implement K-means clustering in C++
        \item Evaluate on standard datasets (Iris, Wine)
        \item Produce simple classifier and outputs for inspection
    \end{itemize}
\end{frame}

\begin{frame}\frametitle{Implementation}
    \begin{itemize}
        \item Language: C++ (modular source files)
        \item Key modules: \texttt{CentroidInitializer}, \texttt{KMeansClassifier}, data IO, \texttt{main}
        \item Design: initialization, iterative update, convergence check
        \item Parallelized compute with MPI for distance and reduction operations
    \end{itemize}
\end{frame}

\begin{frame}\frametitle{Datasets}
    \begin{itemize}
        \item Iris: 150 samples, 3 classes
        \item Wine: 1599 samples, 11 classes
        \item CSV inputs in the \texttt{data/} folder; outputs saved as CSVs
    \end{itemize}
\end{frame}

\begin{frame}\frametitle{Method}
    \begin{itemize}
        \item Standard K-means with Euclidean distance
        \item Centroids updated until convergence
        \item Alternative initializers available in \texttt{CentroidInitializer}
    \end{itemize}
\end{frame}

\begin{frame}\frametitle{Centroid initializer (K-means++)}
    \begin{itemize}
        \item Implements \texttt{k-means++} : selects initial cluster centroids using sampling based on an empirical probability distribution of the points' contribution to the overall inertia
        \item First centroid chosen uniformly at random from data
        \item For each point, compute squared distance to nearest chosen centroid
        \item Select next centroid with probability proportional to that squared distance
        \item Repeat weighted sampling until \texttt{k} centroids are chosen to speed convergence
    \end{itemize}
\end{frame}

\begin{frame}\frametitle{Evaluation}
    \begin{itemize}
        \item Compare cluster labels with ground truth
        \item Metrics: purity and qualitative inspection of CSV outputs
        \item Output examples: \texttt{iris-clusters.csv}, \texttt{wine-clusters.csv}
    \end{itemize}
\end{frame}

\begin{frame}\frametitle{Silhouette Score: Definition and Implementation}
    \begin{itemize}
        \item Measures clustering quality without ground truth labels
        \item For each point $i$: compute $a(i)$ = mean distance to points in same cluster
        \item Compute $b(i)$ = mean distance to nearest other cluster
        \item Silhouette coefficient: $s(i) = \frac{b(i) - a(i)}{\max(a(i), b(i))}$
        \item Range: $[-1, 1]$; higher is better
        \item Implemented in \texttt{silhouette.cpp}: \texttt{euclidean\_distance}, \texttt{calculate\_point\_silhouette}
        \item Helper function computes per-point coefficient; overall score is average
    \end{itemize}
\end{frame}

\begin{frame}\frametitle{Silhouette Score: Multiple Strategies}
    \begin{itemize}
        \item \textbf{Complete Serial} (\texttt{complete\_serial}): computes all points sequentially
        \item \textbf{Complete Parallel} (\texttt{complete\_parallel}): OpenMP parallelization with \texttt{\#pragma omp parallel for} and dynamic scheduling
        \item \textbf{Approximate Serial} (\texttt{approximate\_serial}): random sampling with \texttt{std::shuffle} for faster evaluation
        \item \textbf{Approximate Parallel} (\texttt{approximate\_parallel}): combines sampling with OpenMP for large datasets
        \item Trade-off: accuracy vs. computational cost
    \end{itemize}
\end{frame}

\begin{frame}\frametitle{Results}
    \begin{itemize}
        \item Iris: clear separation for several runs
        \item Wine: partial separation and sensitivity to initialization
        \item Implementation is pedagogical and functional for small datasets
    \end{itemize}
\end{frame}

\begin{frame}\frametitle{Parallelization details}
    \begin{itemize}
        \item MPI initialized in \texttt{main.cpp} with \texttt{MPI\_Init}, \texttt{MPI\_Comm\_rank}, \texttt{MPI\_Comm\_size}
        \item Rank 0 handles I/O and sequential reference run
        \item Data distributed via \texttt{MPI\_Scatterv}; labels gathered with \texttt{MPI\_Gatherv}
        \item Centroids broadcast with \texttt{MPI\_Bcast}; master computes new centroids and broadcasts updates
        \item Timings logged: sequential and parallel runtimes measured in \texttt{KMeansClassifier::fit}
    \end{itemize}
\end{frame}

\begin{frame}\frametitle{Benchmarks (iris dataset, $k=3$)}
    \begin{itemize}
        \item Measured on \texttt{data/dataset-iris.csv} with 3 clusters (results from repository benchmark)
        \item Table: elapsed wall time, exit code and silhouette score (averaged)
    \end{itemize}
    \vspace{3mm}
    \begin{center}
        \begin{tabular}{r r r r}
            \hline
            \# procs & elapsed s & exit code & silhouette \\
            \hline
            1        & 1.974626  & 0         & 0.552592   \\
            2        & 0.239624  & 0         & 0.552592   \\
            4        & 0.250402  & 0         & 0.262306   \\
            \hline
        \end{tabular}
    \end{center}
        \vspace{2mm}
    \begin{itemize}
        \item Observed speedup: $\sim 8\times$ from 1 to 2 procs; diminishing returns at 4 procs for this small dataset
    \end{itemize}
\end{frame}

\end{document}